\section*{Capítulo \ref{ch:b1:angular-momentum} --- Momento angular}

\begin{solution}{exo:b1:angular-momentum:1}
$Fd = 50 \times 0.30 = \qty{15}{N.m}$; $Fd\sin 60^\circ = \qty{13}{N.m}$;
ao longo do cabo a linha de ação passa pela porca: $0$.
\end{solution}

\begin{solution}{exo:b1:angular-momentum:2}
$L = mrv = 5.97\times10^{24} \times 1.50\times10^{11} \times 2.98\times10^4 =
\qty{2.7e40}{kg.m^2/s}$; $\dd A/\dd t = L/2m = rv/2 = \qty{2.2e15}{m^2/s}$.
\end{solution}

\begin{solution}{exo:b1:angular-momentum:3}
$r = \ell\sin 30^\circ = \qty{0.50}{m}$, $\omega = \qty{3.4}{rad/s}$:
$L_z/m = r^2\omega = \qty{0.84}{m^2/s}$. A linha de ação da tensão
passa pelo pivô (\omterm{def:b1:angular-momentum:moment}{momento} nulo em relação a qualquer eixo que passe por ele); o
peso é paralelo ao eixo vertical (\omterm{def:b1:angular-momentum:moment}{momento} nulo em relação a ele): $L_z$
se conserva.
\end{solution}

\begin{solution}{exo:b1:angular-momentum:4}
$J_1 = 4.0 + 2 \times 3.0 \times 0.64 = \qty{7.8}{kg.m^2}$, $J_2 = 4.0 +
2 \times 3.0 \times 0.04 = \qty{4.2}{kg.m^2}$; $\omega_2 = \omega_1J_1/J_2 =
1.85\omega_1$: \qty{3.7}{} voltas por segundo. $E_k = L^2/2J$: razão
$J_1/J_2 = 1.85$; são os músculos, ao recolher os braços contra a
tendência centrífuga, que realizam o trabalho.
\end{solution}

\begin{solution}{exo:b1:angular-momentum:5}
$L_\Delta = m\ell^2\dot\theta$; tensão: momento nulo (encontra o eixo);
peso: $-mg\ell\sin\theta$ (\omterm{def:b1:angular-momentum:moment}{braço de alavanca} $\ell\sin\theta$, restaurador).
$m\ell^2\ddot\theta = -mg\ell\sin\theta$.
\end{solution}

\begin{solution}{exo:b1:angular-momentum:6}
$r_pv_p = r_av_a$: $v_p/v_a = r_a/r_p = 1.034$; com $(v_p + v_a)/2 \approx
\qty{29.8}{km/s}$: $v_p = \qty{30.3}{km/s}$, $v_a = \qty{29.3}{km/s}$.
\end{solution}

\begin{solution}{exo:b1:angular-momentum:7}
$L = mrv$ se conserva: $v_2 = v_1r_1/r_2 = \qty{4.0}{m/s}$, $\omega_2 = v_2/r_2
= \qty{16}{rad/s}$ (partindo de \qty{4}{rad/s}). $E_1 = \qty{0.40}{J}$, $E_2 =
\qty{1.6}{J}$: a puxada realiza \qty{1.2}{J}. A tensão é radial, mas o
deslocamento do disco passa a ter uma componente radial (ele entra em espiral): o trabalho
$\vect T\cdot\dd\vect\ell$ deixa de ser nulo.
\end{solution}

\begin{solution}{exo:b1:angular-momentum:8}
\omterm{def:b1:angular-momentum:moment}{Momentos} em relação ao pivô (comprimentos medidos a partir dele): criança $30 \times 1.5 =
45$ (à esquerda); peso da prancha no seu centro, $0.5$ à direita: $20 \times
0.5 = 10$ (à direita); segunda criança a $x$ à direita: $45x$. Equilíbrio:
$45 = 10 + 45x$, $x = \qty{0.78}{m}$. O pivô sustenta o peso total,
$95 \times 9.81 = \qty{932}{N}$.
\end{solution}

\begin{solution}{exo:b1:angular-momentum:9}
Massas: $m_1a = m_1g - T_1$, $m_2a = T_2 - m_2g$; polia: $J\dot\omega =
(T_1 - T_2)R$ com $a = R\dot\omega$. Somando: $a = (m_1 - m_2)g/(m_1 + m_2 +
J/R^2) = 9.81/(3.0 + 1.0) = \qty{2.5}{m/s^2}$ (contra \qty{3.3}{m/s^2});
$T_1 = m_1(g - a) = \qty{14.7}{N}$, $T_2 = m_2(g + a) = \qty{12.3}{N}$ ---
desiguais, pois a polia precisa de um \omterm{def:b1:angular-momentum:moment}{momento} resultante.
\end{solution}

\begin{solution}{exo:b1:angular-momentum:10}
$r_1v_1 = r_2v_{\perp 2}$: $v_{\perp 2} = 40 \times 0.5/10 = \qty{2.0}{km/s}$.
A \omterm{cor:b1:angular-momentum:central}{lei das áreas} dá apenas a componente perpendicular; a radial
exige a conservação da energia (\cref{ch:b1:central-forces}).
\end{solution}

\begin{solution}{exo:b1:angular-momentum:11}
$J\omega$ constante com $J \propto R^2$: $T' = T(R'/R)^2 = 25 \times 86400
\times (10^4/7\times10^8)^2 = \qty{4.4e-4}{s}$. Velocidade equatorial $2\pi R'/T'
= \qty{1.4e8}{m/s}$, metade da velocidade da luz: o colapso não pode conservar
todo o \omterm{def:b1:angular-momentum:L}{momento angular} --- boa parte é expelida (e a relatividade entra em cena).
\end{solution}

\begin{solution}{exo:b1:angular-momentum:12}
$L_z = mr^2\omega$, e $r$ varia: $\dd L_z/\dd t = 2mr\dot r\omega \neq 0$.
Isso deve ser o \omterm{def:b1:angular-momentum:moment}{momento} da reação transversal $N$ da haste: $rN =
2mr\dot r\omega$, logo $N = 2m\dot r\omega$ --- a haste empurra a conta lateralmente
enquanto ela desliza para fora (o termo de Coriolis do \cref{ch:b1:non-inertial-frames}).
\end{solution}

\begin{solution}{pb:b1:angular-momentum:1}
\textbf{1.} $m\omega^2r = GMm/r^2$: $\omega = \sqrt{GM/r^3} = \sqrt{3.98\times
10^{14}/5.66\times10^{25}} = \qty{2.65e-6}{rad/s}$, período $2\pi/\omega =
\qty{2.37e6}{s} = \qty{27.4}{d}$.

\textbf{2.} $L_{\text{orb}} = 7.35\times10^{22} \times 1.475\times10^{17}
\times 2.65\times10^{-6} = \qty{2.9e34}{kg.m^2/s}$.

\textbf{3.} $v = \sqrt{GM/r}$, $L = mrv = m\sqrt{GMr}$; $\dd L/\dd r =
\tfrac12 m\sqrt{GM/r} = L/2r$.

\textbf{4.} $J = 0.33 \times 5.97\times10^{24} \times 4.06\times10^{13} =
\qty{8.0e37}{kg.m^2}$; $\Omega = \qty{7.29e-5}{rad/s}$; $L_{\text{rot}} =
\qty{5.8e33}{kg.m^2/s}$, um quinto de $L_{\text{orb}}$.

\textbf{5.} $J_m \approx 0.4 \times 7.35\times10^{22} \times 3.0\times10^{12} =
\qty{8.9e34}{kg.m^2}$, vezes $\omega$: \qty{2.4e29}{kg.m^2/s}, $10^{-5}$ do
valor orbital: desprezível.

\textbf{6.} A atração do Sol age sobre o centro de massa do sistema
(sem \omterm{def:b1:angular-momentum:moment}{momento} em relação a ele, em primeira ordem); o seu efeito de maré sobre o par é
pequeno: o $L$ total se conserva com a precisão necessária aqui.

\textbf{7.} A protuberância mais próxima, adiantada em relação à linha Terra--Lua, é puxada
pela Lua com uma força cuja componente tangencial se opõe à rotação
da Terra.

\textbf{8.} \omterm{def:b1:angular-momentum:moment}{Momento} negativo sobre a rotação: a Terra desacelera e o dia
se alonga.

\textbf{9.} A protuberância puxa a Lua para a frente: \omterm{def:b1:angular-momentum:moment}{momento} positivo sobre a
órbita, e $L_{\text{orb}}$ cresce.

\textbf{10.} $L_{\text{orb}} \propto \sqrt r$ cresce, logo $r$ cresce; $v \propto
1/\sqrt r$ diminui.

\textbf{11.} A puxada para a frente realiza trabalho positivo, elevando a energia
da Lua; uma órbita mais alta é uma órbita mais lenta: a energia vai para a altura
(potencial), mais do que a \omterm{thm:b1:work-and-energy:ke}{energia cinética} perdida.

\textbf{12.} Acoplamento de maré da Terra: dia igual ao mês, sem protuberância
adiantada, sem \omterm{def:b1:angular-momentum:moment}{torque}.

\textbf{13.} $\Delta\Omega/\Omega = -\qty{2.3e-3}{s}/\qty{86164}{s} =
\num{-2.7e-8}$ por século.

\textbf{14.} $\Delta L_{\text{rot}} = L_{\text{rot}}\Delta\Omega/\Omega =
\qty{-1.55e26}{kg.m^2/s}$ por século.

\textbf{15.} $\Delta L_{\text{orb}} = +\qty{1.55e26}{kg.m^2/s}$ por século.

\textbf{16.} $\Delta r = 2r\,\Delta L/L = 2 \times 3.84\times10^8 \times 1.55
\times10^{26}/2.9\times10^{34} = \qty{4.1}{m}$ por século, ou \qty{4.1}{cm/yr}:
próximo dos \qty{3.8}{cm/yr} medidos (o alongamento do dia também tem
uma parte não devida às marés).

\textbf{17.} $\Delta T/T = \tfrac32\Delta r/r = 1.6\times10^{-8}$: $2.37\times
10^6 \times 1.6\times10^{-8} = \qty{0.04}{s}$ por século.

\textbf{18.} $\Delta E_{\text{rot}} = L_{\text{rot}}\Delta\Omega = 5.8\times10^{33}
\times 7.29\times10^{-5} \times (-2.7\times10^{-8}) = \qty{-1.1e22}{J}$ por
século.

\textbf{19.} $\Delta E_{\text{orb}} = GMm\Delta r/2r^2 = 2.93\times10^{37}
\times 4.1/(2 \times 1.475\times10^{17}) = \qty{+4.1e20}{J}$ por século.

\textbf{20.} Saldo $-1.1\times10^{22} + 4\times10^{20} \approx \qty{-1.06e22}{J}$
por século: $3.4\times10^{12}\ \unit{W}$, dissipados como calor pelo atrito
de maré nos oceanos e na Terra sólida.

\textbf{21.} Um quinto da potência da humanidade e um catorze avos do fluxo
geotérmico --- um aquecedor pequeno, mas permanente.

\textbf{22.} $L_{\text{tot}} = 3.5\times10^{34}$. Em $r = 5.5\times10^8$:
$\omega = \sqrt{GM/r^3} = \qty{1.55e-6}{rad/s}$, $J\omega = 1.2\times10^{32}$,
$m\sqrt{GMr} = 7.35\times10^{22} \times 4.68\times10^{11} = 3.44\times10^{34}$;
soma $3.45\times10^{34}$: coerente. Período $2\pi/\omega = \qty{4.1e6}{s}
= \qty{47}{d}$.

\textbf{23.} As marés do Sol continuam freando a Terra abaixo da taxa
orbital da Lua, de modo que o \omterm{def:b1:angular-momentum:moment}{torque} Terra--Lua se inverteria; além disso,
o Sol terá evoluído muito antes (a escala de tempo é de dezenas de
bilhões de anos).

\textbf{24.} A Terra levantou marés na Lua, que frearam a sua rotação
até que uma face ficasse voltada para nós; sendo mais leve e mais próxima de um parceiro
mais pesado, o acoplamento da Lua foi muito mais rápido.

\textbf{25.} Conserva-se: o \omterm{def:b1:angular-momentum:L}{momento angular} total. É trocado: \omterm{def:b1:angular-momentum:L}{momento
angular} da rotação da Terra para a órbita da Lua, cerca de \qty{1.6e26}{SI}
por século. É dissipada: a \omterm{thm:b1:work-and-energy:em}{energia mecânica}, alguns terawatts, como calor
de maré.
\end{solution}
